\section*{Kapitel 1 --- Frames}
\addcontentsline{toc}{section}{Kapitel 1 --- Frames}

\solution{ex:1:1}{Bygg en frame}
\ans{a} \code{LEN = 0}, \code{TYPE = 0x02}, \code{DST = 0x25}, \code{SRC = 0x17},
\code{SEQ = 0x7F05}. Checksumman blir
$\mathtt{A5}+\mathtt{F7}+\mathtt{00}+\mathtt{02}+\mathtt{25}+\mathtt{17}+\mathtt{7F}+\mathtt{05}
= 606 = \mathtt{0x025E}$:

\begin{console}
A5 F7 00 02 25 17 7F 05 02 5E
\end{console}

\ans{b} Adresserna byter plats, \code{LEN = 1}, \code{TYPE = 0x03}, en databyte \code{0x16}.
Checksumman blir $630 = \mathtt{0x0276}$:

\begin{console}
A5 F7 01 03 17 25 7F 05 16 02 76
\end{console}

\ans{c} Tio bytes: de åtta headerbytesen plus två checksummabytes, med \code{LEN = 0}.

\solution{ex:1:2}{Fälten och deras syfte}
Utan \textbf{SOF} går det inte att veta var en frame börjar i strömmen. Utan \textbf{LEN} vet
mottagaren inte när payloaden tar slut, eftersom den läser en byte i taget. Utan \textbf{TYPE} går
det inte att avgöra vad framen betyder. Utan \textbf{DST} vet ingen vem den är till, utan
\textbf{SRC} kan mottagaren inte svara, och utan \textbf{SEQ} går svar inte att para ihop med
fråga och dubbletter inte att känna igen. Utan \textbf{CHK} upptäcks inga överföringsfel.

De två som går att ta bort i ett system med exakt två noder är \textbf{DST} och \textbf{SRC}:
med bara två parter är avsändare och mottagare givna av vem som tog emot. Det är också precis
vad som gör dem oumbärliga så snart en tredje nod tillkommer.

\solution{ex:1:3}{Endianness}
\ans{a} \code{05 7F} --- x86 lagrar little-endian, alltså lägsta byten först, medan protokollet
kräver \code{7F 05}.
\ans{b} Båda sidor har samma fel. Serialiseringen skriver \code{05 7F} och deserialiseringen läser
tillbaka det som \code{0x7F05}, så felen tar ut varandra så länge båda ändarna är byggda ur samma
kod på samma arkitektur.
\ans{c} Så fort den andra sidan är något annat: en annan arkitektur, en annan implementation, eller
en bussanalysator som läser strömmen enligt specifikationen. Då blir \code{0x7F05} till
\code{0x057F}.

\solution{ex:1:4}{Checksummans gränser}
\ans{a} Nej. En summa är oberoende av ordningen, så två ombytta bytes ger samma summa.
\ans{b} Nej. Summan är oförändrad, eftersom det som lades till på ett ställe drogs bort på ett
annat.
\ans{c} Till exempel: byt \code{DST} och \code{SRC} med varandra i framen från övning~1.1, eller
ändra \code{0x25} till \code{0x26} och \code{0x17} till \code{0x16}. Båda ger identisk checksumma.
En CRC skulle upptäcka bådadera, vilket är skälet till att CAN använder en.
